Once the delayed neutron count rate is generated, it is converted into \gls{DNP} group parameters.
Fitting the group parameters requires a non-linear least squares method, as both the group half-lives and yields must be determined.
The equation describing the problem appears in Equation \eqref{eq:delnu-fit}, recreated here:
\begin{equation}
\dot{n}_d(t) \approx \epsilon F_s \sum_{k=1}^{K}  \bar{\nu}_{d, k} \lambda_k e^{-\lambda_k t}.
\nonumber
\end{equation}

An important consideration for this process is the selection of time values for generating the group fits.
Increasing the number of time nodes raises computational costs, but too few nodes result in a poorer fit \cite{parish1999status}.
Section \ref{sec:time-eval} shows the sensitivity study for the time nodes, for both the node spacing and the maximum time to measure delayed neutrons.

Another consideration for the non-linear least squares solve is the covariance, which is larger for the pulse irradiation fit than for the saturation irradiation fit.
This results from multiplying the group yield by the group decay constant.
Attempts to replace this multiplication with an independent term produce poor fits, as the multiplication is not independent and is necessary for accurate modeling.

The non-linear least squares process also involves linearization and normalization of the data.
The normalization is handled by taking the delayed neutron count rate, dividing by the fission term \footnote{The fission term is either the net number of fissions or the fission rate, depending on the irradiation type.}, and taking the natural logarithm.
The normalized count rate is then passed into the non-linear least squares fitting.
The linearization is handled by taking the natural logarithm of the delayed neutron count rate.
This treatment reduces both the covariance and the $\chi^2$ of the fit.
Bounds from 0 to 1000 are set for each term to speed up the process, while convergence tolerances are set slightly above machine tolerance at $2.23 \times 10^{16}$ to ensure a well-refined solution.
The convergence criteria include changes in the independent variables and the norm of the gradient, with termination occurring when either criterion meets the set tolerance.


Sample self-multiplication, mentioned in Section \ref{sec:measurementofdnps}, is another consideration.
However, since the fission counts are taken directly from OpenMC tallies, self-multiplication is included and does not introduce error.

The final consideration in \gls{DNP} group reconstruction is uncertainty.
This work tracks uncertainties using the uncertainties package \cite{uncertaintiespackage}.
These uncertainties include uncertainties in cross sections, fission yields, half-lives, and emission probabilities.
These uncertainties are used in the $\chi^2$ minimization, as shown in Equation \eqref{eq:minthis}.


%\begin{description}
%    \item $t_f$ is the final time after irradiation, starting at $t=0$ immediately following irradiation
%    \item $b_t$ is the delayed neutron count vector at the row where the time value is $t$
%    \item $A_t$ is the coefficient matrix at the row where the time value is $t$
%    \item $x_t$ is the solution vector at the row where the time value is $t$
%\end{description}

Once the group fits and their uncertainties are generated, the final step to obtain the finalized group parameters is to combine the saturation and pulse irradiation results.
Saturation irradiation provides better statistics for longer-lived groups, while pulse irradiation is better for shorter-lived groups.
Brady combines the results by using four long-period groups from saturation irradiation and four short-period groups from pulse irradiation.
The groups are renormalized at the second group, where the two longest-lived group half-lives and yield ratios use the saturation results \cite{brady1988evaluation}.
The remaining groups use the pulse irradiation results.

This work uses a similar but slightly simplified approach.
The group parameters are first calculated for both pulse and saturation irradiation.
Then, the two longest-lived groups use the saturation irradiation results, while the four shortest-lived groups use the pulse irradiation results.
